% Part: turing-machines
% Chapter: machines-computations
% Section: representing-tms

\documentclass[../../../include/open-logic-section]{subfiles}

\begin{document}

\olfileid{tur}{mac}{rep}
\olsection{نمایش ماشین‌های تورینگ}

\begin{explain}
ماشین‌های تورینگ را می‌توان به‌صورت دیداری با \emph{نمودارهای حالت}
نمایش داد. این نمودارها از خانه‌های حالت ساخته شده‌اند که پیکان‌ها
آن‌ها را به هم وصل می‌کنند. چنان‌که انتظار می‌رود، هر خانهٔ حالت
نمایندهٔ حالتی از ماشین است. هر پیکان نمایندهٔ دستوری است که از آن
حالت قابل اجراست و جزئیات دستور بالا یا پایین پیکان مناسب نوشته
می‌شود. ماشین زیر را در نظر بگیرید که تنها دو حالت درونی $q_0$ و
$q_1$ و یک دستور دارد:
\[
\begin{tikzpicture}[->,>=stealth',shorten >=1pt,auto,node distance=2.8cm,
                    semithick]
  \tikzstyle{every state}=[fill=none,draw=black,text=black]

  \node[initial,state]   (A)              {$q_0$};
  \node[state]   (B) [right of=A] {$q_1$};

  \path (A) edge  node {\TMtrans{\TMblank}{\TMstroke}{\TMright}} (B);
\end{tikzpicture}
\]
به یاد آورید ماشین تورینگ سر خواندن--نوشتن و نواری دارد که ورودی روی
آن نوشته شده است. دستور را می‌توان چنین خواند: \emph{اگر در حالت
$q_0$ یک~$\TMblank$ می‌خوانی، یک~$\TMstroke$ بنویس، به راست برو و
به حالت $q_1$ وارد شو}. این هم‌ارز آن است که تابع گذار
$\tuple{q_0,\TMblank}$ را به $\tuple{q_1,\TMstroke,\TMright}$ بفرستد.
\end{explain}

\begin{ex}
\emph{ماشین زوج}: ماشین تورینگ زیر متوقف می‌شود اگر و تنها اگر شمار
$\TMstroke$های روی نوار زوج باشد (با این فرض که همهٔ $\TMstroke$ها
پیش از نخستین $\TMblank$ روی نوار آمده‌اند).
\[
\begin{tikzpicture}[->,>=stealth',shorten >=1pt,auto,node distance=2.8cm,
                    semithick]
  \tikzstyle{every state}=[fill=none,draw=black,text=black]

  \node[initial,state]         (A)              {$q_0$};
  \node[state]         (B) [right of=A] {$q_1$};

  \path (A) edge [bend left] node {\TMtrans{\TMstroke}{\TMstroke}{\TMright}} (B)
        (B) edge [loop above] node {\TMtrans{\TMblank}{\TMblank}{\TMright}} (B)
            edge [bend left] node {\TMtrans{\TMstroke}{\TMstroke}{\TMright}} (A);
\end{tikzpicture}
\]

نمودار حالت با تابع گذار زیر متناظر است:
\begin{align*}
\delta(q_0, \TMstroke) & = \tuple{q_1, \TMstroke, \TMright},\\
\delta(q_1, \TMstroke) & = \tuple{q_0, \TMstroke, \TMright},\\
\delta(q_1, \TMblank)  & = \tuple{q_1, \TMblank, \TMright}
\end{align*}
\end{ex}

\begin{explain}
ماشین بالا تنها هنگامی متوقف می‌شود که ورودی شمار زوجی از خط‌ها باشد.
وگرنه ماشین (از نظر نظری) تا ابد به کار ادامه می‌دهد. برای هر ماشین و
ورودی می‌توان \emph{پیکربندی‌های} ماشین را گام‌به‌گام دنبال کرد تا
خروجی تعیین شود. بعداً تعریفی رسمی از پیکربندی‌ها می‌دهیم. فعلاً
می‌توان به‌طور شهودی پیکربندی‌ها را دنباله‌ای از نمودارها دانست که
حالت ماشین را در هر لحظهٔ کار نشان می‌دهند. پیکربندی‌ها محتوای نوار،
حالت ماشین و جای سر خواندن--نوشتن را نشان می‌دهند.

پیکربندی‌های ماشین زوج را هنگامی دنبال کنیم که با ورودی چهار
$\TMstroke$ آغاز می‌شود. در این حالت انتظار داریم ماشین متوقف شود.
سپس ماشین را با ورودی سه $\TMstroke$ اجرا می‌کنیم که در آن تا ابد
کار خواهد کرد.

ماشین در حالت~$q_0$ آغاز می‌کند و چپ‌ترین~$\TMstroke$ را می‌خواند.
حالت آغازین ماشین را می‌توان چنین نمایش داد:
\[
\TMendtape \TMstroke_0 \TMstroke \TMstroke \TMstroke \TMblank \ldots
\]
پیکربندی بالا روشن است. ماشین در حالت صفر آغاز می‌کند و چپ‌ترین
$\TMstroke$ را می‌خواند. این با زیرنویس نام حالت روی نخستین
$\TMstroke$ نمایش داده شده است. دستور قابل اجرا در این نقطه
$\delta(q_0,\TMstroke)=\tuple{q_1,\TMstroke,\TMright}$ است؛ پس ماشین
روی نوار به راست می‌رود و به حالت~$q_1$ تغییر می‌کند.
\[
\TMendtape \TMstroke \TMstroke_1 \TMstroke \TMstroke \TMblank \ldots
\]
چون ماشین اکنون در حالت~$q_1$ یک~$\TMstroke$ می‌خواند، باید دستور
$\delta(q_1,\TMstroke)=\tuple{q_0,\TMstroke,\TMright}$ را «دنبال»
کنیم. حاصل پیکربندی زیر است:
\[
\TMendtape \TMstroke \TMstroke \TMstroke_0 \TMstroke \TMblank \ldots
\]
با ادامهٔ کار ماشین، قواعد دوباره به همان ترتیب اعمال می‌شوند و دو
پیکربندی زیر را به دست می‌دهند:
\[
\TMendtape \TMstroke \TMstroke \TMstroke \TMstroke_1 \TMblank \ldots
\]
\[
\TMendtape \TMstroke \TMstroke \TMstroke \TMstroke \TMblank_0 \ldots
\]
ماشین اکنون در حالت~$q_0$ یک~$\TMblank$ می‌خواند. از نمودار گذار
به‌آسانی می‌بینیم هیچ دستوری برای اجرا نیست و بنابراین ماشین متوقف
شده است. یعنی ورودی پذیرفته شده است.

اکنون فرض کنید ماشین را با ورودی سه $\TMstroke$ آغاز کنیم. چند
پیکربندی نخست مشابه‌اند، زیرا همان دستورها با تفاوت اندکی در ورودی
نوار اجرا می‌شوند:
\[
\TMendtape \TMstroke_0 \TMstroke \TMstroke \TMblank \ldots
\]
\[
\TMendtape \TMstroke \TMstroke_1 \TMstroke \TMblank \ldots
\]
\[
\TMendtape \TMstroke \TMstroke \TMstroke_0 \TMblank \ldots
\]
\[
\TMendtape \TMstroke \TMstroke \TMstroke \TMblank_1 \ldots
\]
ماشین اکنون از همهٔ $\TMstroke$ها گذشته و در حالت~$q_1$ یک
$\TMblank$ می‌خواند. چنان‌که نمودار نشان می‌دهد، دستوری به شکل
$\delta(q_1,\TMblank)=\tuple{q_1,\TMblank,\TMright}$ وجود دارد. چون
نوار در سمت راست تا بی‌نهایت با $\TMblank$ پر است، ماشین اجرای این
دستور را \emph{برای همیشه} ادامه می‌دهد؛ در حالت~$q_1$ می‌ماند و
هرچه بیشتر به راست می‌رود. ماشین هرگز متوقف نمی‌شود و ورودی را
نمی‌پذیرد.
\end{explain}

\begin{explain}
مهم است توجه کنیم که همهٔ ماشین‌ها متوقف نمی‌شوند. اگر توقف یعنی تمام
شدن دستورهای قابل اجرای ماشین، می‌توانیم صرفاً با تضمین وجود پیکانی
خروجی برای هر نماد در هر حالت، ماشینی بسازیم که هرگز متوقف نشود.
می‌توان ماشین زوج را با افزودن دستوری برای خواندن $\TMblank$ در~$q_0$
چنان تغییر داد که بی‌پایان کار کند.
\end{explain}

\begin{ex}
\[
\begin{tikzpicture}[->,>=stealth',shorten >=1pt,auto,node distance=2.8cm,
                    semithick]
  \tikzstyle{every state}=[fill=none,draw=black,text=black]

  \node[initial,state] (A)              {$q_0$};
  \node[state]         (B) [right of=A] {$q_1$};

  \path
  (A) edge [bend left] node {\TMtrans{\TMstroke}{\TMstroke}{\TMright}} (B)
      edge [loop above] node {\TMtrans{\TMblank}{\TMblank}{\TMright}} (A)
  (B) edge [loop above] node {\TMtrans{\TMblank}{\TMblank}{\TMright}} (B)
      edge [bend left] node {\TMtrans{\TMstroke}{\TMstroke}{\TMright}} (A);
\end{tikzpicture}
\]
\end{ex}

\begin{explain}
جدول‌های ماشین راه دیگری برای نمایش ماشین‌های تورینگ‌اند. در جدول
ماشین، الفبای نوار روی محور~$x$ و مجموعهٔ حالت‌های ماشین روی محور~$y$
نمایش داده می‌شود. درون جدول و در تقاطع هر حالت و نماد، باقی دستور---
حالت تازه، نماد تازه و جهت حرکت---نوشته می‌شود. جدول ماشین تعیین حالت
و نمادی را که ماشین در آن متوقف می‌شود آسان می‌کند. هر جای خالی جدول
نقطه‌ای ممکن برای توقف ماشین است. برخلاف نمودارهای حالت و مجموعه‌های
دستور، که نقاط توقف ماشین در آن‌ها همیشه فوراً آشکار نیست، همهٔ نقاط
توقف با یافتن جاهای خالی جدول ماشین به‌سرعت شناسایی می‌شوند.
\end{explain}

\begin{ex}
جدول ماشین برای ماشین زوج چنین است:
\[
\centering
\begin{tabular}{lllll}
\cline{2-4}
\multicolumn{1}{l|}{}      & \multicolumn{1}{c|}{$\TMblank$}                
& \multicolumn{1}{c|}{$\TMstroke$}     & \multicolumn{1}{c|}{$\TMendtape$}   &  \\ \cline{1-4}
\multicolumn{1}{|l|}{$q_0$} & \multicolumn{1}{l|}{}                          
& \multicolumn{1}{l|}{\TMtrans{\TMstroke}{q_1}{\TMright}} & 
\multicolumn{1}{l|}{\phantom{\TMtrans{\TMstroke}{q_1}{\TMright}}} &  \\ \cline{1-4}
\multicolumn{1}{|l|}{$q_1$} & \multicolumn{1}{l|}{\TMtrans{\TMblank}{q_1}{\TMright}} 
& \multicolumn{1}{l|}{\TMtrans{\TMstroke}{q_0}{\TMright}} & 
\multicolumn{1}{l|}{\phantom{\TMtrans{\TMstroke}{q_1}{\TMright}}} &  \\ \cline{1-4}
\end{tabular}
\]
چنان‌که می‌بینیم، ماشین هنگام خواندن~$\TMblank$ در حالت~$q_0$ متوقف
می‌شود.
\end{ex}

\begin{explain}
تا اینجا تنها ماشین‌هایی را در نظر گرفته‌ایم که ورودی را می‌خوانند و
می‌پذیرند. بااین‌حال، ماشین‌های تورینگ هم توان خواندن و هم نوشتن
دارند. نمونهٔ چنین ماشینی (در میان نمونه‌های بسیار) \emph{دوبرابرکن}
است. دوبرابرکن هنگامی که با بلوکی از $n$ عدد $\TMstroke$ روی نوار
آغاز کند، بلوکی از $2n$ عدد $\TMstroke$ بیرون می‌دهد.
\end{explain}

\begin{ex}
  \ollabel{ex:doubler} پیش از ساخت ماشین دوبرابرکن، باید
  \emph{راهبردی} برای حل مسئله بیابیم. چون ماشین (چنان‌که صورت‌بندی
  کرده‌ایم) نمی‌تواند به یاد بسپارد چند $\TMstroke$ خوانده است، به
  راهی برای پیگیری همهٔ $\TMstroke$های نوار نیاز داریم. یک راه این
  است که خروجی را با یک~$\TMblank$ از ورودی جدا کنیم. آنگاه ماشین
  می‌تواند نخستین $\TMstroke$ ورودی را پاک کند، از باقی ورودی بگذرد،
  یک $\TMblank$ باقی بگذارد و دو $\TMstroke$ تازه بنویسد. سپس
  بازمی‌گردد، دومین $\TMstroke$ ورودی را می‌یابد و آن را نیز دوبرابر
  می‌کند. به ازای هر یک $\TMstroke$ ورودی، دو $\TMstroke$ خروجی
  می‌نویسد. با پاک‌کردن ورودی در حین کار تضمین می‌کنیم هیچ
  $\TMstroke$ای از قلم نیفتد یا دوبار دوبرابر نشود. پس از پاک‌شدن کل
  ورودی، $2n$ عدد $\TMstroke$ روی نوار می‌ماند. نمودار حالت ماشین
  تورینگ حاصل در \olref{fig:doubler} آمده است.
  \begin{figure}
\[
\begin{tikzpicture}[->,>=stealth',shorten >=1pt,auto,node distance=2.8cm,
                    semithick]
  \tikzstyle{every state}=[fill=none,draw=black,text=black]

  \node[initial,state] (1)              {$q_0$};
  \node[state]         (2) [right of=1] {$q_1$};
  \node[state]         (3) [right of=2] {$q_2$};
  \node[state]         (4) [below of=3] {$q_3$};
  \node[state]         (5) [left of=4]  {$q_4$};
  \node[state]         (6) [left of=5]  {$q_5$};

  \path (1) edge node {\TMtrans{\TMstroke}{\TMblank}{\TMright}} (2)
    (2) edge [loop above] node {\TMtrans{\TMstroke}{\TMstroke}{\TMright}} (2)
      edge node {\TMtrans{\TMblank}{\TMblank}{\TMright}} (3)
    (3) edge [loop above] node {\TMtrans{\TMstroke}{\TMstroke}{\TMright}} (3)
        edge node {\TMtrans{\TMblank}{\TMstroke}{\TMright}} (4)
    (4) edge [loop below] node {\TMtrans{\TMblank}{\TMstroke}{\TMleft}} (4)
        edge node {\TMtrans{\TMstroke}{\TMstroke}{\TMleft}} (5)
    (5) edge [loop below]  node {\TMtrans{\TMstroke}{\TMstroke}{\TMleft}} (5)
        edge              node {\TMtrans{\TMblank}{\TMblank}{\TMleft}} (6)
    (6) edge [loop below] node {\TMtrans{\TMstroke}{\TMstroke}{\TMleft}} (6)
        edge              node {\TMtrans{\TMblank}{\TMblank}{\TMright}} (1);
\end{tikzpicture}
\]
\caption{ماشین دوبرابرکن}
\ollabel{fig:doubler}
\end{figure}
\end{ex}

\begin{prob}
ورودی دلخواهی انتخاب کنید و پیکربندی‌های ماشین دوبرابرکن در
\olref[tur][mac][rep]{ex:doubler} را گام‌به‌گام دنبال کنید.
\end{prob}

\begin{prob}
ماشین تورینگی با الفبای $\{\TMendtape,\TMblank,A,B\}$ طراحی کنید که
هر رشته از $A$ها و $B$ها را که در آن شمار $A$ها با شمار $B$ها برابر
است \emph{و} همهٔ $A$ها پیش از همهٔ $B$ها آمده‌اند بپذیرد (یعنی روی
آن متوقف شود)، و هر رشته‌ای را که شمار $A$هایش با شمار $B$ها برابر
نیست یا $A$ها پیش از همهٔ $B$ها نیامده‌اند رد کند (یعنی روی آن متوقف
نشود). (مثلاً ماشین باید $AABB$ و $AAABBB$ را بپذیرد، اما $AAB$ و
$AABBAABB$ را رد کند.)
\end{prob}

\begin{prob}
ماشین تورینگی با الفبای $\{\TMendtape,\TMblank,A,B\}$ طراحی کنید که
هر رشتهٔ $\alpha$ از $A$ها و $B$ها را ورودی بگیرد و آن را تکثیر کند
تا خروجی به شکل $\alpha\alpha$ تولید شود. (مثلاً ورودی $ABBA$ باید
خروجی $ABBAABBA$ بدهد.)
\end{prob}

\begin{prob}
\emph{به‌ترتیب الفبا؟} ماشین تورینگی با الفبای
$\{\TMendtape,\TMblank,A,B\}$ طراحی کنید که با گرفتن دنباله‌ای متناهی
از $A$ها و $B$ها بررسی کند آیا همهٔ $A$ها در چپ همهٔ $B$ها قرار
دارند یا نه. ماشین باید رشتهٔ ورودی را روی نوار باقی بگذارد و اگر
رشته «به‌ترتیب الفبا» است متوقف شود، وگرنه تا ابد در حلقه بماند.
\end{prob}

\begin{prob}
\emph{مرتب‌کنندهٔ الفبایی:} ماشین تورینگی با الفبای
$\{\TMendtape,\TMblank,A,B\}$ طراحی کنید که دنباله‌ای متناهی از $A$ها
و $B$ها را ورودی بگیرد و آن‌ها را چنان بازآرایی کند که همهٔ $A$ها در
چپ همهٔ $B$ها باشند. (مثلاً دنبالهٔ $BABAA$ باید به $AAABB$ و دنبالهٔ
$ABBABB$ به $AABBBB$ تبدیل شود.)
\end{prob}

\end{document}
